% ID: SEL-D0029-OA-E04
% ESTADO: cotejada. Segunda pasada visual completa y borrador compilado revisados.
% Fuente: Selectividad/matematicas_aplicadas/suplente septiembre/Examen-Matemáticas Aplicadas a las Ciencias Sociales.pdf
% SHA-256: 4b8884e307df4f5ddaf079e0f10be3d3f74d4793874eb03b355d34f4856579cb
% Página: 1; opción A; ejercicio original 4.

La calificación que obtiene el alumnado en una determinada asignatura sigue una distribución Normal de media $\mu$ y desviación típica 3 puntos.
\begin{enumerate}[label=\alph*)]
\item \textbf{(1.5 puntos)} Se toma una muestra aleatoria simple de 100 alumnos, resultando una calificación media de 5.7 puntos. Calcule un intervalo de confianza para estimar $\mu$ a un nivel de confianza del 95\%.
\item \textbf{(1 punto)} Determine el tamaño mínimo que debe tener una muestra aleatoria para poder estimar $\mu$ con un error máximo de 0.5 puntos y un nivel de confianza del 99\%.
\end{enumerate}
